\chapter{Физическое обоснование абсолютных условий}
\label{ch:axiom-rationale}

\section{Зачем аксиомы, если есть закон \texorpdfstring{$g$}{g}}

Модель строится в два слоя.
На \emph{микро}уровне $\Mlayer$ --- счётный носитель состояний и локальный, обратимый шаг $g$.
На \emph{macro}уровне $\Tlayer$ --- грубый readout (усреднение, binomial blur, Madelung-представление),
из которого восстанавливаются поля, волновые уравнения и константы SI.

Закон $g$ не выбирается «из головы'' как ansatz cellular automaton.
Он --- \emph{замыкание} системы ограничений, которые уже заданы известной физикой:
каузальность и локальность (СТО), сохранение информации (унитарность QM),
термодинамика (второй и третий законы), квантование заряда и спина,
дискретная аналитичность (голоморфность комплексного поля),
симметрии PCT и U(1)$_{\mathrm{vac}}$, геометрический принцип наименьшего действия.

\paragraph{Абсолютные условия A1--A16.}
Это не «настройки симулятора'' и не fitted parameters.
Каждое условие --- формализация физического требования к \emph{допустимому} локальному закону.
Если условие нарушено, модель теряет соответствующий класс явлений на $\Tlayer$
(скорость света, унитарность, необратимость entropy, заряд $e$, spin-$1/2$, и т.д.).

\paragraph{Порядок изложения.}
Настоящая глава --- \emph{мотивация}: откуда в физике берётся каждая группа условий.
Следующая глава (\cref{ch:axioms}) --- формальный реестр A1--A16 с уравнениями.
Геометрия носителя, теорема дискретности и явный вид $g$ --- в главах~\ref{ch:carrier}--\ref{ch:evolution}.

\section{Дискретный субстрат вместо континуума на \texorpdfstring{$\Mlayer$}{M}}

\begin{remark}[Информационный предел]
Комбинация каузальности, конечной энергии в объёме и квантовой неопределённости
накладывает верхнюю границу плотности информации (Bekenstein bound, holographic scaling).
На планковских масштабах трактовать состояние как непрерывное поле $z(x)\in\mathbb{C}$ на $x\in\mathbb{R}^3$
--- значит postulate бесконечный объём фазового пространства на точку,
несовместимый с конечным числом distinguishable конфигураций в causally connected объёме.
\end{remark}

На $\Mlayer$ состояние --- \emph{конечный} регистр на ячейке $\PlanckCell$ (алфавит $Z_N[i]$),
эволюция --- перестановка счётного множества конфигураций, не поток в континууме.
Непрерывные поля, дифференциальные уравнения и «предел $\hl\to 0$'' --- только readout $\Tlayer$
на \emph{фиксированной} решётке; см. \cref{cor:no-continuum} в гл.~\ref{ch:foundations}.

\paragraph{Следствие для выбора условий.}
Любой допустимый $g$ должен:
\begin{itemize}
\item быть \textbf{локальным} и \textbf{каузальным} (информация не опережает light-cone);
\item \textbf{сохранять} норму состояния (нет dissipation информации на micro-шаге);
\item иметь \textbf{конечный} алфавит и \textbf{обратимый} шаг (нет необратимого стирания bit);
\item уважать \textbf{термодинамические} запреты (нет абсолютного нуля, есть потолок плотности);
\item воспроизводить \textbf{топологию} заряда и \textbf{спин} $1/2$ на одной ячейке.
\end{itemize}
Именно эти требования группируются в A1--A16.

\section{Каузальность, локальность, изотropия}
\label{sec:rat-causality}

\paragraph{Специальная теория относительности.}
Скорость распространения сигнала ограничена $c$.
На дискретном носителе это означает: за один такт $\htick$ возмущение не может пройти
дальше одного шага $\hl$ по решётке; вторая координационная оболочка за тик запрещена
(иначе «диагональ за тик'' --- superluminal stencil).

\paragraph{Локальность поля.}
Continuum-уравнения (Klein--Gordon, Dirac, гидродинамика) в пределе $\Tlayer$
восстанавливаются из \emph{локальных} stencils: новое значение в точке $x$ зависит
только от окрестности $N(x)$, не от всего пространства.

\paragraph{Изотropия macro-света.}
Micro-шаг может быть анизотропным на решётке; macro-readout усредняет направления
до изотropного «света'' с $\kgeom=1/\sqrt{2}$ на FCC (см. \cref{ch:carrier,ch:macro}).

\begin{table}[ht]
\centering
\caption{Каузальность и геометрия $\Rightarrow$ условия}
\label{tab:causality-axioms}
\begin{tabular}{@{}ll@{}}
\toprule
Физическое требование & Условие \\
\midrule
Light-cone: один тик $\le$ один hop $\hl$; $|N|=12$ в $(3{+}1)$ & A1 \\
Локальный stencil, не global oracle & A2 \\
Macro-Lorentz / изотropный readout & A12 \\
\bottomrule
\end{tabular}
\end{table}

\section{Квантовая механика как сохранение информации}
\label{sec:rat-qm}

\paragraph{Унитарность.}
В QM эволюция сохраняет $\braket{\psi|\psi}$.
На $\Mlayer$ это не postulate «волновой функции'', а \emph{инвариант} $N(t)=\sum_x z^\dagger z$:
линейная часть $g$ --- product локальных unitaries; нелинейность --- только фаза $e^{i\Phi(z)}$,
не damping $|z|$.

\paragraph{Спинор $\mathbb{C}^2$, не скаляр.}
Одна комплексная фаза $z\in\mathbb{C}$ не различает оборот $2\pi$ и тождество для fermion:
нужен спинор $z\in\mathbb{C}^2$ и SU(2)-gate $R(\Phi)$ (A4, A16).
U(1)$_{\mathrm{charge}}$ --- глобальная фаза; локальное столкновение --- SU(2).

\paragraph{Обратимость micro-шага.}
Унитарность на конечном алфавите требует явного $g^{-1}$.
Leapfrog второго порядка на $\mathbb{Z}$ с kick-ledger --- минимальная форма,
совместимая с A3 и дискретным временем (A13).

\paragraph{Детерминизм и «квантовая случайность''.}
На $\Mlayer$ нет RNG: $g$ --- детерминированная функция (следствие A2).
Born и интерференция на $\Tlayer$ --- статистика грубого прибора и хаотической фазовой динамики (A5, A6),
не фундаментальные кости в $g$ (\cref{prop:determinism}).

\begin{table}[ht]
\centering
\caption{QM $\Rightarrow$ условия}
\label{tab:qm-axioms}
\begin{tabular}{@{}ll@{}}
\toprule
Физическое требование & Условие \\
\midrule
Сохранение $N=\sum|z|^2$; unitary linear + phase nonlinear & A3 \\
Спинор, U(1)$_{\mathrm{charge}}$, SU(2) collision & A4 \\
Обратимый шаг; leapfrog 2-го порядка & A13 \\
Spin-$1/2$: $2\pi\to -1$, Pauli на $\PlanckCell$; $s_0=\hbar/2$ & A16 \\
\bottomrule
\end{tabular}
\end{table}

\section{Термодинамика на ячейке}
\label{sec:rat-thermo}

\paragraph{Третий закон.}
Абсолютный ноль $|z|=0$ недостижим: иначе фазовый gate $\Phi$ теряет смысл ($z e^{i\Phi}=0$ --- deadlock),
и «вакуум'' становится статическим нулём, а не физическим фоном с $\rho_{\mathrm{field}}>0$.
A5 формализует: минимум «потенциала'' не в $z=0$; при $|z|\to 0$ kick $\Phi\to\infty$ --- вечное micro-кипение.

\paragraph{Второй закон (локально).}
Энтропия не убывает на $\varepsilon$-окрестности при \emph{фиксированной} норме $N$.
На $\Mlayer$ это mixing фаз, не dissipation $|z|$ --- иначе A3 нарушен (A6).

\paragraph{Планковский потолок плотности.}
Энергия на ячейку ограничена $u_P$; saturating gate с знаменателем $+\varepsilon$ (A7).
При больших $|z|$ нелинейность затухает --- macro-поле линейно (A8), как в классическом пределе QM.

\begin{table}[ht]
\centering
\caption{Термодинамика $\Rightarrow$ условия}
\label{tab:thermo-axioms}
\begin{tabular}{@{}ll@{}}
\toprule
Физическое требование & Условие \\
\midrule
3-й закон; вакуум = кипящий фон, не $z\equiv 0$ & A5 \\
2-й закон локально; mixing без $N\to 0$ & A6 \\
Потолок $\rho_E\le u_P$; saturating & A7 \\
Macro-линейность при больших $|z|$ & A8 \\
\bottomrule
\end{tabular}
\end{table}

\paragraph{Тепловая смерть $\Mlayer$.}
Комбинация A3, A5, A6, A13 влечёт: $N(t)=\mathrm{const}$, нет shutdown, нет устойчивого $z\equiv 0$
(\cref{sec:theorem-23}). «Тепловая смерть'' наблюдаемого --- phenomena $\Tlayer$ (D3), не смерть micro-субстрата.

\section{Аналитичность, заряд, солитоны}
\label{sec:rat-analytic}

\paragraph{Дискретная голоморфность.}
В continuum holomorphic branch выбирается условиями Коши--Римана.
На решётке $\Arg(\sum z)-\Arg z$ даёт ложные скачки $2\pi$; физически значима holonomy
$\zeta=(\sum_N z)\cdot z^*$ (A9).
CR-единственность на связном $\Lambda$ задаёт глобальное поле --- риманова поверхность $\Lambda\times S^1$.

\paragraph{Квантование заряда.}
Устойчивый vortex --- полюс на $\partial(\hv)$; winding
$n=\frac{1}{2\pi}\oint d\arg z\in\mathbb{Z}$ (A10).
Дробный заряд --- не устойчивое состояние того же $g$ без изменения топологии.

\paragraph{Anti-smear / устойчивость солитона.}
Размазанное (не голоморфное) состояние не устойчиво: gate с $K_P$ и $\Delta\phi_{\min}$
выталкивает полюс --- самофокусировка, не бесконечный smear (A11).

\begin{table}[ht]
\centering
\caption{Аналитичность и топология $\Rightarrow$ условия}
\label{tab:analytic-axioms}
\begin{tabular}{@{}ll@{}}
\toprule
Физическое требование & Условие \\
\midrule
Holonomy $\zeta$; дискретная CR; глобальное поле & A9 \\
Заряд $n\in\mathbb{Z}$ на $\partial(\hv)$ & A10 \\
Устойчивость vortex; запрет smear & A11 \\
\bottomrule
\end{tabular}
\end{table}

\section{Симметрии и геометрия coupling'ов}
\label{sec:rat-sym}

\paragraph{P, C, T, U(1)$_{\mathrm{vac}}$.}
Стандартные дискретные симметрии должны быть symmetries \emph{закона} $g$,
не только Lagrangian $\Tlayer$-readout (A14).

\paragraph{Принцип наименьшего действия / без fitted knobs.}
Безразмерные константы ($\kgeom$, $\gamma$, $\alpha^*$) --- следствия геометрии FCC-носителя
и bit-budget ячейки $B_{\hv}=2\pi/\ln 2$, не свободные параметры fit (A15).
Подробный вывод $\alpha$, масс и SI --- \cref{ch:si-sm}.

\begin{table}[ht]
\centering
\caption{Симметрии и действие $\Rightarrow$ условия}
\label{tab:sym-axioms}
\begin{tabular}{@{}ll@{}}
\toprule
Физическое требование & Условие \\
\midrule
P, C, T, U(1)$_{\mathrm{vac}}$ на $g$ & A14 \\
$\kgeom$, $\gamma=1/|N|$, $\alpha^*=1+1/(4\pi)$ из геометрии & A15 \\
\bottomrule
\end{tabular}
\end{table}

\section{Сводная карта: физика $\to$ реестр}

\begin{table}[ht]
\centering
\caption{Полная карта мотивации (формулировки --- \cref{ch:axioms})}
\label{tab:full-map}
\small
\begin{tabular}{@{}clp{7.2cm}@{}}
\toprule
 & Условие & Физический корень \\
\midrule
A1 & Каузальность & Light-cone; один $\htick$ --- один hop \\
A2 & Локальность & Local stencil; нет action-at-a-distance \\
A3 & Унитарность & Сохранение информации; $N=\mathrm{const}$ \\
A4 & U(1)/SU(2) & Заряд и спин на $\mathbb{C}^2$ \\
A5 & 3-й закон & Нет $|z|=0$; вакуумное кипение \\
A6 & 2-й закон & Mixing энтропии при fixed $N$ \\
A7 & Потолок $u_P$ & Planck density bound \\
A8 & Macro-линейность & Классический предел \\
A9 & CR / holonomy & Голоморфность; $\zeta$ вместо $\atan2$ \\
A10 & Заряд $n\in\mathbb{Z}$ & Winding на $\partial(\hv)$ \\
A11 & Солитон & Anti-smear; $K_P$ \\
A12 & Изотropия & Macro-Lorentz readout \\
A13 & Обратимость & $g^{-1}$; leapfrog \\
A14 & PCT, U(1)$_{\mathrm{vac}}$ & Дискретные symmetries \\
A15 & Действие & Геометрия $\Rightarrow$ couplings \\
A16 & Fermion & $2\pi\to -1$; Pauli; $s_0=\hbar/2$ \\
\bottomrule
\end{tabular}
\end{table}

\section{Что условия \emph{не} делают}

\begin{remark}[Замыкание vs ansatz]
A1--A16 \emph{не} задают явную таблицу lookup для $g$.
Они задают \emph{класс} допустимых $g$.
Единственный закон, удовлетворяющий всем условиям на $Z_{512}[i]$, --- замыкание
leapfrog + holonomy kick + $R(\Phi)$ (\cref{thm:g-law}).
\end{remark}

\begin{remark}[Носитель и шкала]
Выбор FCC, $|N|=12$, bit-budget $B_{\hv}$ и идентификация $\hl$ с $\ell_P$ ---
геометрия (\cref{ch:carrier}) и слой SI (\cref{ch:si-sm}), не отдельные «аксиомы произвола''.
Теорема дискретности (\cref{thm:discreteness}) показывает согласованность
носителя с A1, A2, A10, A16.
\end{remark}

\begin{remark}[Meta vs model]
Вопросы «замка наблюдателя'', UI-кадра, космологического «первого тика'' наблюдаемой Вселенной
--- epistemic/meta-слой; они \emph{не} добавляют параметров в $g$
и не входят в A1--A16.
Micro-пена A5 (вечные тики) $\neq$ начало наблюдаемой космологии.
\end{remark}

\paragraph{Дальше.}
Формальные определения, уравнения и доказательства теоремы о невозможности тепловой смерти $\Mlayer$
--- \cref{ch:axioms}.
